\chapter{Evals}
\label{ch:evals}

\section*{What this chapter is for}

Evaluation is the discriminator. \reported{Portfolio evidence of an evaluation
harness is used by some employers specifically to separate people who have
built something from people who have followed a
tutorial}{field-guide-get-hired}, and \reported{not starting a take-home with
evaluation is reported as a red flag}{field-guide-take-home}. It is also
\reported{the steepest part of the curve for engineers arriving from backend
work.}{field-guide-backend}

\section*{The questions}

\begin{bankq}{What is an eval, and how is it different from a test?}
\qasked
The opening question, and it sounds like a definition request. It is not.

\qtested
Whether you understand what non-determinism does to the instruments you
already have.

\qstaff
A test asserts that a deterministic thing produced a specific result, and a
failure means something changed. Neither half survives here. The same input
produces different output, so an assertion on the output is either too loose to
catch anything or too tight to stay green.

What replaces it is a shift in what you assert about. Assert on \emph{structure
and behaviour} rather than on text: which tools were called, in what order,
with which arguments, what was emitted, what the outcome was. Those are
stable across rewordings and they are the things that matter.

Where the output itself is the product and must be judged, you move to a
measured quantity rather than a boolean: a score over a set, compared against a
recorded baseline, where the question is not ``did it pass'' but ``is it worse
than it was''.

\judgement{The sentence that makes this land in an interview: a test asks
whether the code is correct, and an eval asks whether the system is still as
good as it was. The second is a comparison, so it is meaningless without a
baseline \dash{} and that is why the first thing an evaluation suite needs is
not cases but a recorded previous state.}

\qsenior
``Evals are tests for non-deterministic systems, usually with a model grading
the output.'' Not wrong, and it goes straight to the expensive, unstable
mechanism while skipping the cheap deterministic one that should gate every
change.

\qcontested
How much can be made deterministic is genuinely disputed and depends on the
product. \judgement{Most systems have far more deterministic surface than their
teams use, because the visible part is the prose and the prose is what gets
graded.}

\qdrill
For a model-driven feature you know, list five things you could assert
deterministically about one interaction. If you cannot reach five, the system
is not telling you enough about what it did.
\end{bankq}

\begin{bankq}{How do you know your eval suite actually works?}
\qasked
The follow-up, and the best question in this chapter.

\qtested
Whether you have applied any scepticism to your own instrument. Most people
have not, and this is the question where a prepared candidate is unmistakable.

\qstaff
\textbf{A suite that has never failed is a suite nobody has tested.} The way to
know is to break the system deliberately and confirm the suite notices.

Build variants with one thing wrong each: one that skips the confirmation, one
that ignores a permission, one that leaks an internal identifier. Run the suite
against each. A variant that survives has found a gap in your coverage, and
that is the point of the exercise rather than an embarrassment.

Two details that make it real rather than ceremonial. The broken variant must
not announce itself \dash{} it differs only in the behaviour under test, or you
are checking that you can detect a label. And you must confirm the case passes
against the \emph{correct} system first, or a case that was broken some other
way looks like a successful catch.

\judgement{If you have done this and one of your variants survived, say so. It
is the strongest sentence available on this subject, because it is the sentence
of somebody who tested an instrument and reported what it did rather than what
they hoped.}

The general form is worth stating too: a check nobody has watched fail is not
known to be checking anything. It applies well beyond this domain, and it is
usually the most transferable thing a backend engineer brings.

\qsenior
Describing coverage \dash{} how many cases, how many categories, how broad. It
answers how much the suite looks at, not whether looking works.

\qcontested
How many broken variants is enough. \judgement{Nobody knows. Three is enough to
have the conversation; the number matters far less than having done it once.}

\qdrill
Take any test suite you own and break the thing it tests, in one line, without
changing the test. Time how long until something goes red. If nothing does, you
have learnt something about a suite you trusted this morning.
\end{bankq}

\begin{bankq}{You want a model to grade output. How do you keep that trustworthy?}
\qasked
About model-based grading, and usually sceptical in tone.

\qtested
Whether you treat the grader as another component with its own error rate, or
as an oracle.

\qstaff
Anchor the rubric. ``Rate this one to ten'' produces a number with nothing to
regress against, because nobody \dash{} including the grader \dash{} can say
what six means. Every level has to name a behaviour a second reader could agree
or disagree with. That is what makes the next step possible.

Calibrate before it gates anything. Label a sample by hand, compare with the
grader's labels, and compute agreement. Until agreement is established the
grader reports and does not block. \judgement{And the labels have to be yours:
a grader calibrated against labels produced by another model has been checked
against itself.}

Pin and record. The grader's model, its rubric and its prompt all go into the
report \dash{} a hash is enough \dash{} so a run graded under different
conditions is visible rather than inferred. Two scores from different rubric
versions are not comparable and nothing will tell you.

And structurally: the graded layer is slow, costs money and needs a credential,
so it cannot gate every change. The deterministic layer does that. This layer
runs on a schedule and answers a different question.\source{hamel-llm-judge}

\qsenior
A grader with a good prompt and a sensible scale, spot-checked occasionally. It
is the common implementation, and the phrase ``spot-checked occasionally'' is
where it stops being measurement.

\qcontested
Which agreement statistic, and what threshold, are genuinely contested and the
choice is less important than reporting one. \judgement{More contested, and
more interesting: whether model graders should gate at all. The defensible
position is that they inform and the deterministic layer gates \dash{} a red
build that nobody can reproduce teaches people to re-run it.}

\qdrill
Write one rubric level \dash{} say, a three out of five \dash{} for a quality
you care about, such that two colleagues would assign it to the same outputs.
It is harder than it looks and the difficulty is the finding.
\end{bankq}

\begin{bankq}{What do you assert, beyond the right answer?}
\qasked
Sometimes asked directly, more often reached by following the previous answer.

\qtested
Whether your idea of correct includes things that must not happen.

\qstaff
Roughly a fifth of assertions in a serious suite should be about absence: a
tool that must not be called, an event that must not fire, an internal
identifier that must not appear in user-facing output.

The reason is specific. Consider a case where the system should decline. Assert
only that it declined and you have tested half of it \dash{} an agent that
refuses politely and calls the tool anyway passes. The pair \dash{} refusal
issued, \emph{and} the call absent \dash{} is the whole assertion, and one
without the other is a check with a hole in it shaped exactly like the failure
you were worried about.

\judgement{Absence assertions are also the cheapest thing in this chapter to
add to an existing suite and the thing most likely to catch something on the
day you add them.}

Two more that pay for themselves: that the loop terminated by decision rather
than by hitting its iteration cap, and that every identifier used in a write
appeared in an earlier result rather than being constructed.

\qsenior
A thorough account of asserting correct outputs across categories. Complete on
one side.

\qcontested
Little, on absence. What is contested is how much of this belongs in evaluation
rather than in the system \dash{} and \judgement{the answer is both: the tool
boundary enforces it, and the suite proves the boundary works.}

\qdrill
Take three cases from any suite you have and add the absence half to each. At
least one will fail immediately, and that is the chapter's argument in one
exercise.
\end{bankq}

\kitfig{two-layers}{% Chapter 32. The two layers answer different questions and only one of them
% can gate a change.
\begin{tikzpicture}[node distance=6mm]
  \node[kitbox, minimum width=42mm, minimum height=13mm] (trace)
    {\textbf{One captured trace}\\{\scriptsize tool calls \dash{} arguments \dash{} events \dash{} outcome}};

  \node[kitbox, minimum width=44mm, minimum height=16mm, below left=10mm and -22mm of trace] (det)
    {\textbf{Deterministic}\\{\scriptsize structure and absence}\\{\scriptsize no model, no credential}};
  \node[kitsoft, minimum width=44mm, minimum height=16mm, below right=10mm and -22mm of trace] (judge)
    {\textbf{Judged}\\{\scriptsize anchored rubric, calibrated}\\{\scriptsize pinned and hashed}};

  \draw[kitarrow] (trace) -- (det);
  \draw[kitarrow] (trace) -- (judge);

  \node[kitbox, draw=kitink, minimum width=44mm, below=6mm of det]
    (gate) {\textbf{Gates every change}};
  \node[kitsoft, minimum width=44mm, below=6mm of judge]
    (rep) {\textbf{Reports on a schedule}};

  \draw[kitarrow] (det) -- (gate);
  \draw[kitarrow, draw=kitgrey] (judge) -- (rep);

  \node[kitlabel, anchor=north, text width=46mm] at ($(gate.south)+(0,-1mm)$)
    {\emph{is it still doing the right things?}};
  \node[kitlabel, anchor=north, text width=46mm] at ($(rep.south)+(0,-1mm)$)
    {\emph{is it still as good as it was?}};
\end{tikzpicture}
}{Two layers over one captured trace. Only the deterministic one can gate a change, because only it runs without a model, a credential or a network.}

\section*{What this chapter does not claim}

The reports about evaluation's role as a discriminator and as the steepest part
of the transition come from one corpus, and the calibration practices are
consistent with published practitioner writing.\source{hamel-evals-faq}
\source{anthropic-agent-evals}

The proportion of absence assertions, the position that graders should not
gate, and the instruction to keep calibration labels human are the author's
judgement. The first is a rule of thumb and not a measurement.
